% ------------------------------------------------------------
% Page 76
% ------------------------------------------------------------

\begin{center}
    {\Large\bfseries Cabrillac}

    \vspace{0.15cm}

    {\large Un hameau sur l’Aigoual}
\end{center}

\vspace{0.45cm}

Cabrillac est le rassemblement de quelques maisons à 1200m d’altitude, écrit Yves Grellier,
sur un sol formant replat, sur le versant nord de l’Aigoual, à quelque sept km. de son sommet
par la route. Il est immédiatement dominé par la draille qui monte droit le long de la
montagne vers le Sud Est. Le village surplombe la haute vallée de la Jonte qui traverse et
arrose de nombreux prés….

\vspace{0.3cm}

Cabrillac est un carrefour sur ce col battu des vents, entre la Jonte et le Tarnon, entre
l’Aigoual et le Causse, entre granite, schiste et calcaire, depuis longtemps lieu de passage des
hommes et des bêtes… Une grande foire annuelle s’y tenait le 20 Mai, foire aux chèvres
comme son nom l’indique, et aux cochons. Les gens du Vigan montent à Cabrillac acheter
leur cochon de l’année qu’ils tueront à Noël…. *

\vspace{0.3cm}

C’est en 1827 que René de Bernis neveu du cardinal-ministre vend ses terres de Cabrillac à
trois chefs de famille du lieu : Avesque, Poujol et Verdier : familles encore bien représentées
dans le pays, inaugurant ce que l’on a pu appeler l’âge d’or de Cabrillac, mais qui ne dura
qu’une soixantaine d’années.

Pendant trois quarts de siècle de 1870 à 1950 qui aura vu trois guerres ! Cabrillac s’est vidé de
ses travailleurs, de ses enfants, son terroir s’est rétréci, ses troupeaux amenuisés : des vaches
remplacent, en partie, les si nombreuses brebis.

\vspace{0.3cm}

On peut retenir deux dates particulières pour Cabrillac : \textbf{1953} c’est l’arrivée des premiers
estivants urbains qui ont, pour la plupart, une parenté même lointaine avec les Avesque,
Poujol ou Verdier. Ils achètent et retapent de vieilles écuries ou bergeries inoccupées. Une
population nouvelle va remplacer l’ancienne au moins pendant les mois d’été puis de plus en
plus longtemps avec les retraités et l’allongement du temps de la vie.

\textbf{1980} c’est le décès du dernier paysan.

\vspace{0.3cm}

L’hiver Cabrillac est désormais désert.

\vspace{0.5cm}

\begin{center}
\setlength{\fboxsep}{8pt}
\fbox{%
\parbox{0.88\textwidth}{%
\small
\textbf{*} Mme. Poujol dont les petits enfants habitent Montcamp, Meyrueis, Sainte Enimie, etc. préparait la soupe pour
les fermiers du Causse Méjean qui venaient faire leur provision de bois dans les forêts de l’Aigoual. Ils
arrivaient de Nivoliers ou de la Bégude, de Fretma ou de Gally, d’Aures ou de l’Hom et du Veygaliers…. Les
attelages de bœufs traversaient Cabrillac au petit jour. Sans bruit ils ouvraient la porte Poujol (on ne fermait pas
les portes en ces temps là) et déposaient dans le tiroir de la grande table un morceau de lard pour assaisonner la
soupe des hommes et leur préparer par Mme. Poujol : celle-ci avait coutume de raconter « au morceau de lard je
sais quel est le fermier qui va revenir ». Etait-ce à la quantité, à la fraîcheur du morceau, à la façon dont il
avait été conservé ? Nous ne savons…
}%
}
\end{center}
